\section{結言}
今回の観測窓では、生成AIの能力向上そのもの以上に、その競争を誰がどのように統治するかが中心的な論点になった。Anthropicを離れた研究者の警告、同社自身によるサイバー事案assessment、OpenAIの安全ガバナンス人事が同日に重なり、個人の危機感、企業の自己評価、外部検証の必要性が並行して表れた。

同時に、Muse、Images 2.5、プラグイン、Defense Factoryのようなエージェント・製品展開と、WeatherNext、AlphaGenome、フィンランドの計算基盤投資が進んでいる。生成AIをめぐる競争は、モデル性能だけでなく、実利用、安全、研究、エネルギーを含む基盤全体へ広がっている。

一方、Navier--Stokes関連主張や一部の安全発言、AlphaGenomeの数値には未確認要素が残る。Daily XではX上で何が話題になったかを記録し、一次資料との厳密な照合や採否は週刊版の工程へ委ねる。
